


\documentclass[11pt]{article}
\usepackage{amsgen,amsmath,amstext,amsbsy,amsopn,amsfonts,amsthm,amscd}
%\usepackage{bbm}
\usepackage[dvips]{graphicx}
\usepackage{enumerate}
\setlength{\textwidth}{16.0cm} \setlength{\textheight}{25cm}
\voffset=-2.50cm \hoffset=-2.00cm \setlength{\topmargin}{0.4cm}

%%%%%%%%%%%%%%%%%%%%%%%%%%%%%%%%%%%%%%%%%%%%%%%%%%%%%%%%%%%%%%

\newtheorem{lem1}{Lema}%[chapter]
\newtheorem{th1}{Teorema}%[chapter]
\newtheorem{prop1}{Proposici\'{o}n}%[chapter]
\newtheorem{defi1}{Definici\'{o}n}%[chapter]
\newtheorem{coro1}{Corolario}%[chapter]
\newtheorem{ob1}{Observaci\'{o}n}%[chapter]
\newtheorem{obs1}{Observaci\'{o}n}%[chapter]
\newtheorem{nota1}{Nota}%[chapter]
\newtheorem{eje1}{Ejemplo}%[chapter]
\newtheorem{ejes1}{Ejemplo}%[chapter]
\newtheorem{ejerc1}{Ejercicio}%[chapter]
\newtheorem{ejercs1}{Ejercicio}%[chapter]
\newenvironment{prop}{\begin{prop1}\rm}{\end{prop1}}
\newenvironment{defi}{\begin{defi1}}{\end{defi1}}
\newenvironment{coro}{\begin{coro1}}{\end{coro1}}
\newenvironment{th2}{\begin{th1}}{\end{th1}}
\newenvironment{nota}{\begin{nota1}\rm}{\end{nota1}}
\newenvironment{lem}{\begin{lem1}}{\end{lem1}}
\newenvironment{ob}{\begin{ob1}\rm}{\end{ob1}}
\newenvironment{obs}{\begin{obs1}\rm}{\end{obs1}}
\newenvironment{eje}{\begin{eje1}\rm}{\end{eje1}}
\newenvironment{ejes}{\begin{ejes1}\rm}{\end{ejes1}}
\newenvironment{ejerc}{\begin{ejerc1}\rm}{\end{ejerc1}}
\newenvironment{ejercs}{\begin{ejercs1}\rm}{\end{ejercs1}}
%%%%%%%%%%%%%%%%%%%%%%%%%%%%%%%%%%%%%%%%%%%%%%%%%%%%%%
%\DeclareSymbolFont{cucu}{U}{bbmss}{m}{n}
\DeclareMathOperator{\CC}{\mathbb C}%{\mathalpha}{cucu}{67}
\DeclareMathOperator{\NN}{\mathbb N}%{\mathalpha}{cucu}{78}
\DeclareMathOperator{\RR}{\mathbb R}%{\mathalpha}{cucu}{82}
\DeclareMathOperator{\QQ}{\mathbb Q}%{\mathalpha}{cucu}{81}
\DeclareMathOperator{\ZZ}{\mathbb Z}%{\mathalpha}{cucu}{90}
%%%%%%%%%%%%%%%%%%%%%%%%%%%%%%%%%%%%%%%%%%%%%%%%%%%%%%
\DeclareMathOperator{\sen}{sen}
\DeclareMathOperator{\arcsen}{arcsen}
\DeclareMathOperator{\senh}{senh} \DeclareMathOperator{\fr}{Frac}
%%%%%%%%%%%%%%%%%%%%%%%%%%%%%%%%%%%%%%%%%%%%%%%%%%%%
\newcommand{\ran}[1]{{\em Im\/}(#1)}
\newcommand{\ca}[1]{{\em card\/}(#1)}
\newcommand{\sop}[1]{{\em supp\/}(#1)}
\newcommand{\tra}[1]{{\em traza\/}(#1)}
\newcommand{\dis}[1]{d (#1)}
\newcommand{\inte}[1]{{\em Int\/}(#1)}
\newcommand{\conv}[1]{{\em conv\/}(#1)}
\newcommand{\too}{\longrightarrow}
\newcommand{\oo}{\overline}
\newcommand{\cq}{\begin{flushright} $\Box$ \end{flushright}}
\newcommand{\cqd}{\, \, \rule{7pt}{7pt}}
\def\limsup{\mathop{\overline{\rm lim}}}
\def\liminf{\mathop{\underline{\rm lim}}}
\renewcommand{\theenumii}{\arabic{enumii}}
\renewcommand{\labelenumii}{\theenumi.\theenumii.}
\renewcommand{\theenumiii}{\arabic{enumiii}}
\renewcommand{\labelenumiii}{\theenumi.\theenumii.\theenumiii}
%%%%%%%%%%%%%%%%%%%%%%%%%%%%%%%%%%%%%%%%%%%%%%%%%%%%%%%%%%%%%%%%%
\begin{document}

\section{UNA COTA PARA LOS CEROS DE LOS POLINOMIOS DE SOBOLEV}

\noindent Producto Sobolev continuo: caso secuencialmente dominado.

\noindent Estimaci\'{o}n de la cota de acotaci\'{o}n de ceros de polinomios ortogonales de Sobolev.

\noindent Consideremos el siguiente producto de Sobolev:
$$
  <p(t),q(t)>=  \int_{a}^{b} p(t)q(t)w_0(t)\; dt +  \int_{a}^{b} p'(t)q'(t)w_1(t)\; dt
$$

\noindent donde $w_0(t),w_1(t)$ son pesos continuos de modo que est\'{a}n sequencialmente dominados en el sentido de que: Existe una constante $C>0$ tal que:
$$
  w_1(t) \leq Cw_0(t) , \qquad t\in [a,b]
$$

\noindent Puesto que $\dfrac{w_1(t)}{w_0(t)}$ es una funci\'{o}n continua en $[a,b]$ consideremos
$$
  \alpha= \max \{ \dfrac{w_1(t)}{w_0(t)}, t\in [a,b]\}.
$$

\noindent Consideremos el operador multiplicaci\'{o}n por $t$
$$
  \mathcal{D}: \mathbb{P}[t] \to \mathbb{P}[t], \quad \mathcal{D}(p(t))=tp(t).
$$

\noindent Veamos que $\mathcal{D}$ es acotado en este caso adem\'{a}s:
$$
  \Vert \mathcal{D} \Vert \leq C= (R^{2}+2\alpha)^{1/2}.
$$

$$
  \Vert \mathcal{D}(p(t))\Vert^{2} = \int_{a}^{b} t^2p^2(t)w_0(t)\; dt +  \int_{a}^{b} (tp'(t))^2w_1(t)\; dt= \int_{a}^{b} t^2p^2(t)w_0(t)\; dt +  \int_{a}^{b} (p(t)+tp'(t))^2w_1(t)\; dt \leq
$$
\noindent utilizando que $(a+b)^2\leq 2(a^2+b^2)$ si $a,b\geq 0$ se tiene que:
$$
  \int_{a}^{b} t^2p^2(t)w_0(t)\; dt+2\int_{a}^{b} p(t)^2w_1(t)\; dt +\int_{a}^{b} +t^2p'(t))^2w_1(t)\; dt \leq
$$
$$
  R^{2} \int_{a}^{b} p^2(t)w_0(t)\; dt +  \int_{a}^{b} (p'(t)^2w_1(t)\; dt+ 2\int_{a}^{b} p(t)^2w_1(t)\; dt \leq
$$
$$
  R^{2} \Vert p(t)\Vert^{2} + 2\alpha \int_{a}^{b} p(t)^2w_0(t)\; dt \leq (R^{2}+2\alpha)\Vert p(t) \Vert^{2}
$$

\noindent Por ejemplo, si $[a,b]=[-1,1]$ y $\alpha =1$ quedar\'{\i}a $C=\sqrt{1+2}=\sqrt{3}$.

\begin{enumerate}
  \item Experimentar con ejemplos sequencialmente dominados $w_0(t)=t^n$ y $w_1(t)=t^{2n}$ con $n$ par en el intervalo $[-1,1]$, se aprecia que la c ota es buena ???
  \item Experimentar con ejemplos que no sean sequencialmente dominados $w_0(t)=t^{2n}$ y $w_1(t)=t^{n}$ con $n$ par en el intervalo $[-1,1]$, se aprecian resultados an\'{a}logos al caso anterior.
  \item Observar que en los casos anteriores todos los ceros están en el círculo cerrado de centro el origen y radio $\sqrt{3}$.
\end{enumerate}
\noindent


\section{Conjetura (reunión 13 de marzo )}

\noindent A la vista de los resultados: consideremos el conjunto de los ceros de los polinomios del polinomio ortonormal (ortogonal, etc.) de orden $n$ de Sobolev:
$$
  Z_n=\{ z\in \CC: P_n(z)=0\}= \{ z \in \CC: \vert \mathbf{D}_n - zI_n\vert=0\}.
$$

\noindent es claro que el conjunto de todos los ceros de los polinomios de Sobolev es
$$
  Z=\cup_{n=1}^{\infty} Z_n
$$

\noindent Por otra parte, si consideramos todos los autovalores de $\mathbf{D}_n$ son los ceros de $P_n(z)$, por lo que denotando el autovalor m\'{a}ximo de $\mathbf{D}_n$ como $\lambda_{max}(\mathbf{D_n})$ y denotando por $z_n$ el cero de $P_n(z)$ tal que
$$
  \vert z_n\vert =\lambda_{max}(\mathbf{D}_n)
$$

\noindent El problema de la acotación de los ceros consiste en el problema de la acotación de los $\{z_n\}_{n=1}^{\infty}$:



\bigskip

\noindent {\bf CONJETURA:} productos de Sobolev soportados en intervalos del tipo
$[a_1,a_2],[a_3,a_4]$ si $R=\max \{\vert a_k \vert : k=1,\dots,4\}$

\begin{enumerate}
  \item
        $$
          \lim_{n\to \infty} \vert \lambda_{max}(\mathbf{D}_n)\vert = R;
        $$
  \item Si $\vert a_{k_0} \vert=R$ entonces
        $$
          \lim_{n\to \infty} z_n = a_{k_0},
        $$
        \noindent es decir, $a_{k_0}$ sería un atractor de los ceros con módulo máximo.

\end{enumerate}

\noindent Nótese que son equivalentes:
\begin{enumerate}
  \item $Z$ es acotado.
  \item La sucesión $(\lambda_{max}(\mathbf{D}_n))_{n=1}^{\infty}$ está acotada
\end{enumerate}

\noindent Por lo tanto, se propone para el estudio de la acotación de los ceros de los polinomios ortogonales de Sobolev la acotación de estos autovalores máximos.

\noindent En el caso de que las matrices $\mathbf{D}_n$ sean diagonalizables, lo que ocurre cuando \underline{todas} las raíces de los polinomios ortogonales son distintos. Esto ocurre en los ejemplos estudiados, pero sabemos que no es cierto en el caso de polinomios ortogonales de Sobolev. Tampoco es cierto en el caso de polinomios soportados en medidas en el plano complejo, pero sí lo es en el caso de medidas soportadas enla recta real, por la simetría de la matriz de Jacobi.


\noindent No es cierto, en general, que
$$
  \Vert \mathbf{D}_n \Vert = \lambda_{max}(\mathbf{D}_n)
$$


\noindent Para ello, basta considera las secciones del shift-right. Por  ejemplo,
la sección de orden dos,
$$
  \mathbf{D}_2 = \begin{pmatrix} 0 & 0 \\ 1 & 0 \end{pmatrix}
$$

\begin{enumerate}
  \item $\Vert \mathbf{D}_2 \Vert =1 $
  \item $\lambda_{max}(\mathbf{D}_2)=0$ (el único autovalor es $0$).
\end{enumerate}



\noindent Sin embargo, en el caso de que todas las raíces de los polinomios ortogonales de Sobolev (caso recta real) fuesen distintos entonces las matrices $\mathbf{D}_n$ serían diagonalizables y en este caso sí sería cierto que
$$
  \Vert \mathbf{D}_n \Vert = \lambda_{max}(\mathbf{D}_n)
$$

\noindent Se puede probar que para una matriz con entradas reales:
$$
  \Vert \mathbf{D}_n \Vert= \lambda_{max}(\mathbf{D}_n^{t}\mathbf{D}_n)
$$

\noindent la matriz $\mathbf{D}_n^{t}\mathbf{D}_n$ es una matriz simétrica, autoadjunta, por lo que todos sus autovalores son reales.
\bigskip

\noindent {\bf PROBLEMA:} ¿Es cierto para los casos de productos de Sobolev soportados en la recta real que los ceros de los polinomios ortogonales son todos distintos?
\newpage

\section{Reunión de 27 de marzo}


\noindent Comprobación de todos los pasos:

\begin{prop} Sean $\mu_1,\mu_2$ medidas con soporte infinito  en $\CC$  y consideremos un producto interno
  $$
    <p(z),q(z)>= \int p(z)\overline{q(z)} d\mu_0+ \int p(z)\overline{q(z)} d\mu_1
  $$

  \noindent Sea $\{P_n(z)\}_{n=0}^{\infty}$ la sucesión de polinomios ortonormales. Entonces


  $$
    Z=
    \{ z \in \CC: \text{ there exists } n\in \NN, \varphi_n(z)=0\}\subset \{ z\in \CC  : \;  \vert z \vert < \Vert \mathcal{D} \Vert \}.
  $$
\end{prop}

\begin{proof} Supongamos que  $z_0\in Z$, es una raíz de un polinomio $P_n(z)=(z-z_0)q_{n-1}(z)$; además   $q_{n-1} \neq 0$ y $q_{n-1}(z)$ es de grado menor o igual que $n-1$.
  $$
    <P_n(z),P_n(z)>=<(z-z_0)q_{n-1}(z),(z-z_0)q_{n-1}(z)>=<zq_{n-1}(z),
  $$



  $$
    \vert z_0 \Vert^2 \Vert q_{n-1} \Vert^2 = \Vert zP_{n-1} \Vert^2 - \Vert P_n(z) \Vert^2 < \Vert z P_{n-1}(z) \Vert^2 = \Vert \mathcal{D}(P_{n-1}(z)) \Vert \leq \Vert \mathcal{D} \Vert^2 \Vert P_{n-1} \Vert^2
  $$

  \noindent Since  $\Vert P_{n-1} \Vert \neq 0$ it follows that
  $$
    \vert z_0  \vert^2 <  \Vert \mathcal{D} \Vert^2
  $$

  \noindent as we required.
\end{proof}

\newpage
\section{Reunión 10 de abril}
\noindent Estudio de los casos

$$
  \Vert p \Vert^2_{\mathbf{S}} = \int_{-1}^{1} p^{2}(t)+ \int_{-1}^{1} t^{2} (p'(t))^{2}\; dt + \vert p'(0)\vert^{2}
$$

\noindent En el artículo J.Rodríguez, ...se afirma que el operador multiplicación por $z$ no es acotado, por lo que
$$
  \lim_{n\to \infty} \Vert \mathbf{D}_n \Vert =\infty
$$


\noindent En los siguientes casos, salen por Rodriguez??, también se puede afirmar lo mismo.

$$
  \Vert p \Vert_{\mathbf{S}} = \int_{-1}^{1} p^{2}(t)\; dt + \vert p'(a)\vert^{2}
$$
$$
  \Vert p \Vert_{\mathbf{S}} = \int_{0}^{1} p^{2}(t)\; dt + \vert p'(a)\vert^{2}
$$

\noindent $\bullet$ {\bf ESTUDIAR EL COMPORTAMIENTO DEPENDIENDO DE LA LOCALIZACIÓN DEL ÁTOMO  $a$ : }
$$
  \lim_{n\to \infty} \rho(\mathbf{D}_n),
$$

\noindent Experimentalmente se observa que en el caso :
$$
  \int_{0}^{1} p^{2}(t)\; dt + \vert p'(0) \vert^{2},
$$
$$
  \lim_{n \to \infty} \rho(\mathbf{D}_n)=0.
$$

\noindent {\bf INTENTAR PROBARLO}



\noindent En algunos casos parece ir al átomo, en otros casos, al punto más alejado del origen?? La simetría parece tener que ver.

\bigskip
\bigskip

\noindent $\bullet$ {\bf MUY INTERESANTE:}

\noindent Estudio de la convergencia del \underline{segundo} mayor   de los valores singulares de $\mathbf{D}_n$, que parece no ser el mismo que el del primer autovalor!!!!.

\noindent Dada la matriz $\mathbf{D}_n^{t}\mathbf{D}_n$ que es simétrica y semidefinida positiva consideramos el segundo mayor valor singular, es decir, si
$$
  0\leq \lambda_{n}^{n} \leq \dots \leq \lambda_{2}^{n} \leq \lambda_{1}^{n}
$$



\noindent son los valores singulares de $\mathbf{D}_n$,

\bigskip

\noindent nótese que $\rho(\mathbf{D}_n)= \lambda_{1}^{n}$.

\bigskip


\noindent $\bullet$ {\bf  ESTUDIO DEL COMPORTAMIENTO ASINTÓTICO de: }
$$
  \lim_{n\to \infty} \lambda_{2}^{n}.
$$

\noindent tiene relación con el soporte de las medidas???


\end{document}
\newpage

\section{CASO DE CIRCUNFERENCIAS}

$$
  <p,q>=\int p(z) \overline{q(z)}d{\bf m} + \int_{\mathbf{S}_r} p'(z) \overline{q'(z)}d{\bf m_{0;r}
$$

$$
  \Vert p(z) \Vert^{2}_{\mathbf{S}} = \Vert zp(z) \Vert_{{\bf m}}^{2} + \Vert p'(z) \Vert^{2}_{{\bf m}_{0;r}}=
$$






$$
  \Vert \mathcal{D}(p(z))\Vert_{\mathbf{S}}^{2} = \Vert \vert zp(z)\Vert_{{\bf m}}^{2} + \Vert (zp(z))' \Vert^{2}_{{\bf m}_{0;r}}=
$$

\noindent Puesto que en la primera integral $\vert z \vert=1$ se tiene que:

$$
  =\Vert  p(z) \Vert_{{\bf m}}
  ^{2} + \Vert p(z)+zp'(z) \Vert^{2}_{{\bf m}_{0;r}}
$$
$$
  \leq \Vert  p(z) \Vert^2_{{\bf m}} + 2 \Vert  p(z) \Vert^2_{{\bf m}_{0;r}} + 2\Vert zp'(z) \Vert^{2}_{{\bf m}_{0;r}} \leq \Vert  p(z) \Vert^2_{{\bf m}} + 2 \Vert  p(z) \Vert^2_{{\bf m}_{0;r}} + 2r^2\Vert p'(z) \Vert^{2}_{{\bf m}_{0;r}} \leq
$$
\noindent Si $C= \max \{1,2r^2\}$ se tiene que
$$
  \leq C(\Vert  p(z) \Vert^2_{{\bf m}}+\Vert p'(z) \Vert^{2}_{{\bf m}_{0;r}})+2\Vert  p(z) \Vert^2_{{\bf m}_{0;r}}= C\Vert p(z) \Vert^{2}_{\mathbf{S}}+ \Vert  p(z) \Vert^2_{{\bf m}_{0;r}}
$$

\noindent Utilizaremos ahora la desigualdad polinomial de Wirtinger para la medida ${\bf m}_{0;r}$ que dice que si $q(z)$ es un polinomio con $q(0)=0$ se tiene que:
$$
  \Vert q(z) \Vert^{2} = \int_{\mathbf{S}_r} \left(\sum_{k=0}^{n} a_kz^k \right)= \sum_{k=0}^{n} r^{2k} \vert a_k \vert^{2} \leq
$$
$$
  \sum_{k=0}^{n} r^{2k}  \vert a_k \vert^{2}\leq \sum_{k=1}^{n} kr^{2k-2}\vert a_k \vert^{2} = \int_{\mathbf{S}_r} \vert q'(z)\vert^{2}d{\bf m}_{0;r}=  \Vert q'(z)\Vert^{2}.
$$


\noindent Luego,
$$
  \vert q(0) \vert
$$


\end{document}
